\documentclass{stex}

\libusepackage{naproche}
\libusepackage{copyright}
\libinput{preamble}

\begin{document}

\usemodule[libraries/meta]{preliminaries.ftl}


\section{Propositions}\label{sec:propositions}

Since \ForTheL is essentially a surface language for (unsorted) first-order logic,
it provides symbolic notations for logical connectives and quantifiers which are
described in this section.


\subsection{Truth}\label{sec:truth}

\begin{sparagraph}[style=symdoc,for=TRUE]
  \begin{signature*}
    $\fakeemph{\TRUE}$ is a formula.
  \end{signature*}

  \begin{axiom*}
    $\TRUE$ is true.
  \end{axiom*}
\end{sparagraph}


\subsection{Falsity}\label{sec:falsity}

\begin{sparagraph}[style=symdoc,for=FALSE]
  \begin{signature*}
    $\fakeemph{\FALSE}$ is a formula.
  \end{signature*}

  \begin{axiom*}
    $\FALSE$ is false.
  \end{axiom*}
\end{sparagraph}


\subsection{Negation}\label{sec:negation}

\begin{sparagraph}[style=symdoc,for=NOT]
  \begin{signature*}
    Let $\varphi$ be a formula.
    $\fakeemph{\NOT \varphi}$ is a formula.
  \end{signature*}

  \begin{axiom*}
    Let $\varphi$ be a formula.
    $\NOT\varphi$ is true iff $\varphi$ is false.
  \end{axiom*}
\end{sparagraph}


\subsection{Conjunction}\label{sec:conjunction}

\begin{sparagraph}[style=symdoc,for=AND]
  \begin{signature*}
    Let $\varphi, \psi$ be formulas.
    $\fakeemph{\varphi \AND \psi}$ is a formula.
  \end{signature*}

  \begin{axiom*}
    Let $\varphi, \psi$ be formulas.
    $\varphi \AND \psi$ is true iff $\varphi$ is true and $\psi$ is true.
  \end{axiom*}
\end{sparagraph}


\subsection{Disjunction}\label{sec:disjunction}

\begin{sparagraph}[style=symdoc,for=OR]
  \begin{signature*}
    Let $\varphi, \psi$ be formulas.
    $\fakeemph{\varphi \OR \psi}$ is a formula.
  \end{signature*}

  \begin{axiom*}
    Let $\varphi, \psi$ be formulas.
    $\varphi \OR \psi$ is true iff $\varphi$ is true or $\psi$ is true.
  \end{axiom*}
\end{sparagraph}


\subsection{Implication}\label{sec:implication}

\begin{sparagraph}[style=symdoc,for=IMPLIES]
  \begin{signature*}
    Let $\varphi, \psi$ be formulas.
    $\fakeemph{\varphi \IMPLIES \psi}$ is a formula.
  \end{signature*}

  \begin{axiom*}
    Let $\varphi, \psi$ be formulas.
    $\varphi \IMPLIES \psi$ is true iff $\varphi$ is true implies that $\psi$ is true.
  \end{axiom*}
\end{sparagraph}


\subsection{Equivalence}\label{sec:equivalence}

\begin{sparagraph}[style=symdoc,for=IFF]
  \begin{signature*}
    Let $\varphi, \psi$ be formulas.
    $\fakeemph{\varphi \IFF \psi}$ is a formula.
  \end{signature*}

  \begin{axiom*}
    Let $\varphi, \psi$ be formulas.
    $\varphi \IFF \psi$ is true iff $\varphi$ is true if and only if $\psi$ is true.
  \end{axiom*}
\end{sparagraph}


\subsection{Universal Quantification}\label{sec:universal-quantification}

\begin{sparagraph}[style=symdoc,for=FORALL]
  \begin{signature*}
    Let $x$ be a variable and $\tau$ be a term and $R$ be a binary relation and $\varphi$ be a formula.
    $\fakeemph{\FORALL} x \mathrel{R} \tau : \varphi$ is a formula.
  \end{signature*}

  \begin{axiom*}
    Let $x$ be a variable and $\tau$ be a term and $R$ be a binary relation and $\varphi$ be a formula.
    $\FORALL x \mathrel{R} \tau : \varphi$ is true iff $\varphi[\tau/x]$ is true for all terms $\sigma$ such that $\sigma \mathrel{R} \tau$.
  \end{axiom*}
\end{sparagraph}


\subsection{Existential Quantification}\label{sec:existential-quantification}

\begin{sparagraph}[style=symdoc,for=EXISTS]
  \begin{signature*}
    Let $x$ be a variable and $\tau$ be a term and $R$ be a binary relation and $\varphi$ be a formula.
    $\fakeemph{\EXISTS} x \mathrel{R} \tau : \varphi$ is a formula.
  \end{signature*}

  \begin{axiom*}
    Let $x$ be a variable and $\tau$ be a term and $R$ be a binary relation and $\varphi$ be a formula.
    $\EXISTS x \mathrel{R} \tau : \varphi$ is true iff $\varphi[\tau/x]$ is true for some term $\sigma$ such that $\sigma \mathrel{R} \tau$.
  \end{axiom*}
\end{sparagraph}


\ifinputref\else\printlicense[CcByNcSa]{Marcel Schütz (2026)}\fi

\end{document}
